\chapter{Nasazení}
Tato kapitola se zabývá procesem zprovoznění a nasazení implementované aplikace. Nejprve bude podrobněji představen způsob využití technologie Docker. Následně bude uveden podrobný popis postupu pro spuštění systému v lokálním vývojovém prostředí a bude přiblíženo nastavení integrace s externími službami. V závěru kapitoly bude popsán proces nasazení aplikace s využitím cloudové platformy Render.

\section{Docker}
Celá aplikace je navržena pro provoz v izolovaném prostředí pomocí technologie Docker. Pro optimalizaci procesu sestavení je využíván mechanismus vícefázového sestavení. Tento přístup přináší následující klíčové benefity:

\begin{itemize}
    \item \textbf{Minimalizace času sestavení}: Díky rozdělení procesu do jednotlivých fází je umožněno využití již vytvořených vrstev. Tím je zamezeno opakování redundantních operací, což vede k výrazné časové úspoře při opětovném sestavování.
    \item \textbf{Snížení velikosti výsledného obrazu}: Do finálního obrazu se přenášejí pouze nezbytné soubory a zkompilovaný kód. Jelikož vývojové nástroje a meziprodukty zůstávají v předchozích fázích, je výsledná velikost obrazu minimalizována.
    \item \textbf{Oddělení prostředí}: Definice konkrétních cílů umožňuje z jediného konfiguračního souboru generovat specifické obrazy optimalizované pro odlišné účely.
\end{itemize}

\newpage
Pro účely této práce byly definovány dva cíle sestavení aplikace. Pro potřeby lokálního vývoje byl navržen cíl \texttt{development} a pro produkční prostředí byl specifikován cíl \texttt{production}.

\begin{itemize}
    \item \texttt{development}: Tato konfigurace je rozšířena o nástroj xDebug, který slouží k pokročilému ladění zdrojového kódu v jazyce PHP. Dále je integrováno prostředí Node.js verze 24 pro práci se závislostmi v jazyce JavaScript.
    \item \texttt{production}: Tento cíl byl nastaven jako výchozí a výsledný obraz je optimalizován s důrazem na minimalizaci jeho celkové velikosti. I z tohoto důvodu nebyly do obrazu zahrnuty vývojové nástroje, jako je například prostředí Node.js či adresář \texttt{node\_modules}.
\end{itemize}

Implementace cíle \texttt{production} je součástí výpisu kódu \ref{listing:dockerfile}. Proces sestavení začíná zkopírováním PHP závislostí a sestavených assetů z předchozích fází \texttt{vendor} a \texttt{node}. Následně jsou zkopírovány samotné zdrojové kódy aplikace a je provedeno nastavení vlastnických oprávnění. Příkazem \texttt{CMD} je definována sekvence kroků prováděných při startu kontejneru: nejprve jsou spuštěny databázové migrace, poté je vykonán příkaz \texttt{php artisan optimize} pro optimalizaci výkonu frameworku a nakonec je spuštěn webový server Apache.

\begin{listing}[h!]
    \caption{Cíl production z Dockerfile souboru}\label{listing:dockerfile}
    \begin{minted}{dockerfile}
FROM base AS production
WORKDIR /var/www/html

COPY --from=vendor /var/www/html/vendor/ ./vendor/
COPY --from=node /app/public/build/ ./public/build/
COPY . .

RUN chown -R www-data:www-data bootstrap/cache storage
RUN chmod -R 775 bootstrap/cache storage

CMD composer dump-autoload --optimize --no-dev --strict-ambiguous \
    && php artisan migrate --force \
    && php artisan optimize \
    && exec apache2-foreground
    \end{minted}
\end{listing}

\newpage
\section{Spuštění na lokálním prostředí}
Pro zprovoznění aplikace na lokálním prostředí je nutné provést kroky popsané v této sekci. Předpokladem pro úspěšné spuštění je nainstalovaný nástroj Docker a stažený repozitář se zdrojovým kódem projektu. Veškeré kroky by měly být prováděny v kořenovém adresáři projektu.

\begin{enumerate}
    \item vytvoření konfiguračního souboru \texttt{.env}, obsah lze zkopírovat ze vzorového souboru \texttt{.env.example}

    \item sestavení a spuštění aplikačních kontejnerů a databáze
    \begin{minted}{bash}
docker-compose up -d
    \end{minted}

    \item instalace závislostí pro backendovou část aplikace
    \begin{minted}{bash}
docker-compose exec php php composer install
    \end{minted}

    \item instalace závislostí pro JavaScript a sestavení single page aplikace
    \begin{minted}{bash}
docker-compose exec php npm install
docker-compose exec php npm run build
    \end{minted}

    \item vygenerování bezpečnostního klíče aplikace
    \begin{minted}{bash}
docker-compose exec php php artisan key:generate
    \end{minted}

    \item provedení databázových migrací pro vytvoření schématu
    \begin{minted}{bash}
docker-compose exec php php artisan migrate
    \end{minted}

    V rámci procesu migrace je v systému automaticky vytvořen administrátorský účet s následujícími přihlašovacími údaji:
    \vspace{-2mm}
    \begin{itemize}
        \item e-mail: \texttt{admin@devpulse.com}
        \item heslo: \texttt{password}
    \end{itemize}

    \item naplnění databáze testovacími daty (volitelné)
    \begin{minted}{bash}
docker-compose exec php php artisan db:seed 
    \end{minted}
\end{enumerate}

Po úspěšném provedení výše uvedených kroků je systém dostupný prostřednictvím webového prohlížeče na adrese \texttt{http://localhost/}. V případě potřeby lze všechny běžící kontejnery zastavit příkazem \texttt{docker-compose down}.

\section{Nastavení integrace s externími systémy}
Pro zajištění korektní synchronizace dat je nezbytné provést konfiguraci propojení s externími systémy pro správu verzí. Následující sekce popisují proces získání požadovaných autentizačních údajů a postup jejich nastavení v rámci vytvořené aplikace.

\subsection{GitHub App}
Integrace s platformou GitHub je realizována prostřednictvím registrace nové GitHub App, která může být vytvořena pod osobním účtem či v rámci organizace. V průběhu vytváření aplikace je vyžadováno vyplnění pouze jejího názvu a URL domovské stránky. Pro zajištění funkčnosti synchronizace je následně nezbytné udělit oprávnění v sekci \texttt{Repository permissions}, konkrétně pro~\texttt{Pull requests} v režimu \texttt{Read-only}.

Po dokončení registrace je aplikaci automaticky přidělen unikátní identifikátor \texttt{Client ID}. Pro získání privátního klíče je nutné provést jeho manuální vygenerování v administračním rozhraní aplikace. Tyto údaje je následně nezbytné vložit do~konfiguračního souboru \texttt{.env} pod klíči \texttt{GITHUB\_APP\_CLIENT\_ID} a \texttt{GITHUB\_APP\_PRIVATE\_KEY}. Závěrečným krokem je instalace aplikace do cílového účtu či organizace na platformě GitHub. Tímto úkonem je získán identifikátor \texttt{installation\_id}, který je nutné uvést při vytváření nové instance systému správy verzí v administraci vytvořené aplikace.

\subsection{GitLab}
V případě platformy GitLab je proces autentizace a následný přístup k datům realizován prostřednictvím přístupových tokenů. Z hlediska architektury vytvořeného systému je doporučeno využívat osobní či skupinové tokeny. V případě využití projektových tokenů by byla vyžadována konfigurace nové instance pro správu verzí pro každý sledovaný repozitář.

Při generování přístupového tokenu je nezbytné udělit oprávnění \texttt{read\_api}, které je vyžadováno pro synchronizaci dat. Na rozdíl od platformy GitHub není vygenerovaný token ukládán do konfiguračního souboru \texttt{.env}, ale je vkládán přímo do příslušného formuláře v průběhu vytváření nové instance systému pro správu verzí softwaru.

\subsection{E-mailový server}
Nedílnou součástí konfigurace produkčního prostředí je nastavení e-mailového serveru. Ve výchozím nastavení, které je definováno v souboru \texttt{.env.example}, je pro odesílání e-mailů nastaven ovladač \texttt{log}. Při využití tohoto ovladače avšak nedochází k reálné distribuci zpráv příjemcům, ale k jejich zápisu do~souboru \texttt{storage/logs/laravel.log}. Pro reálné odesílání e-mailů je nutné zvolit například ovladač \texttt{smtp}. V konfiguračním souboru \texttt{.env} musí být v tomto případě definovány hodnoty proměnných, které specifikují hostitele, port a autentizační údaje k příslušnému poštovnímu serveru. Seznam požadovaných proměnných je uveden na následujícím výpisu kódu.

\begin{minted}{bash}
MAIL_MAILER=smtp
MAIL_HOST=
MAIL_PORT=
MAIL_USERNAME=
MAIL_PASSWORD=
MAIL_FROM_ADDRESS=
\end{minted}

\section{Render}
Pro hostování aplikace byla zvolena platforma Render. Jedná se o cloudovou službu typu Platform as a Service, která umožňuje provozovat webové aplikace a související služby bez nutnosti správy vlastní serverové infrastruktury. V rámci realizovaného řešení je na této platformě hostována nejen samotná aplikační vrstva, ale také relační databáze MariaDB a plánovač úloh, který zajišťuje pravidelnou synchronizaci dat ze systémů pro správu verzí softwaru.

Celý proces nasazení byl plně automatizován díky přímé integraci s verzovacím systémem GitHub. Konfigurace zajišťuje, že při přidání commitu do~větve \texttt{master} dojde k automatickému sestavení a nasazení nové verze aplikace. Pro~korektní běh systému bylo nutné na straně poskytovatele definovat proměnné prostředí, které odpovídají obsahu lokálního souboru \texttt{.env}.

Z bezpečnostního hlediska je produkční verze, na rozdíl od lokálního vývojového prostředí, provozována výhradně na šifrovaném protokolu HTTPS. Výsledná aplikace byla zpřístupněna na doméně \texttt{devpulse.paukertlukas.cz}. Pro odesílání e-mailových notifikací byl systém propojen s externím SMTP serverem služby Mailtrap.
